\chapter{Fazit}\label{chap:Schluss}
Im Rahmen dieser Arbeit wurde eine \glsentrylong{por}-Schaltung für einen neu entwickelten Chip entworfen, der sowohl analoge als auch digitale Funktionsblöcke enthält. Aufgrund der unterschiedlichen Anforderungen dieser beiden Bereiche wurden separate \gls{por}-Schaltungen entwickelt und gezielt an die jeweiligen Versorgungsspannungen angepasst. Der Fokus lag dabei auf der Verzögerung und einer sofortigen Freigabe des Reset-Signals. Die ursprüngliche widerstandsbasierte Referenzspannung erwies sich in der Corner-Simulation als zu empfindlich gegenüber Prozess- und Temperaturschwankungen. Durch die Umstellung auf eine transistorbasierte Referenzspannung mit nativen Transistor konnte die Stabilität der Schaltung deutlich verbessert werden. Die umfangreichen Corner- und \glsentrylong{mc}-Simulationen bestätigten, dass die entworfenen \gls{por}-Schaltungen die Spezifikationen auch unter variierenden \gls{pvt}-Bedingungen erfüllen. Damit wurde ein wichtiger und grundlegender Beitrag zur sicheren Initialisierung sowohl der analogen als auch der digitalen Teil des Chips geleistet.

Vor der Freigabe für das Layout muss jedoch die Auffälligkeit aus der \gls{mc}-Process-Simulation untersucht werden. Diese zeigte unerwartet Abweichungen der Schwellwerte, welche vermutlich durch den nativen Transistor, der zur Erzeugung der Referenzspannung eingesetzt wird. Eine genaue Analyse und Korrektur sind erforderlich, bevor die Schaltung in das Layout integriert werden kann. Zudem ist es essenziell, die \gls{por}-Schaltungen im Layout-Prozess sorgfältig zu integrieren und dabei mögliche Fehlerquellen zu minimieren. Insbesondere parasitäre\footnote{\url{https://resources.pcb.cadence.com/blog/2019-how-parasitic-capacitance-and-inductance-affect-your-signals} --\,zuletzt am 27.~Februar~2025 erfolgreich aufgerufen.} Widerstände und Kapazitäten, Kopplungen zu benachbarten Signalleitungen oder unzureichende Abschirmungen können die Funktion der Schaltung erheblich beeinflussen. Daher sind weitere Simulationen erforderlich, um die Stabilität auch im finalen Layout sicherzustellen.

Trotz der noch bestehenden Herausforderungen konnte im Rahmen dieser Arbeit ein wesentlicher Beitrag zur Entwicklung des neuen Chips geleistet werden. Die grundlegende Funktionalität der \gls{por}-Schaltungen wurde nachgewiesen, und die Schaltungen erfüllen weitestgehend die Anforderungen. Damit ist eine zuverlässige Spannungsüberwachung sichergestellt, sodass der Chip fehlerfrei in Betrieb genommen werden kann. 
%
% Das folgende Label ist für die korrekte Zählung der Seitenzahl im
% Autorenreferat zwingend an dieser Stelle erforderlich.
% Stellen Sie sicher, dass es am Ende dieses Kapitels steht!
\label{contentPages:End}